\chapter{Response to Phase-2 Examiner Feedback}
\label{app:examiner-response}

\noindent
\begin{tabular}{@{} >{\bfseries}p{0.28\textwidth} p{0.68\textwidth} @{}}
Name of the Candidate & Rajat Sharma \\
Roll Number & CH24M559 \\
Title of the Thesis & Estimating Remaining Useful Life and State of Health of
EV Batteries by Hybrid Physics-Informed Deep Learning Framework \\
Degree & Master of Technology in Industrial Artificial Intelligence \\
\end{tabular}

\section*{Abstract (of the final thesis)}

This thesis develops a hybrid physics-informed framework that regularizes a
lightweight Gated Recurrent Unit (GRU) with empirical degradation physics for
SOH and RUL estimation of lithium-ion cells. Two physics mechanisms are
implemented and compared: a fixed closed-form prior (route~A, using the
integrated Verhulst and double-exponential laws) and a learnable ODE residual
(route~B). Ten models are evaluated on the NASA PCoE and CALCE datasets under
a leak-free leave-one-cell-out protocol with symmetric inputs and a
matched-capacity control. The learnable route-B hybrid attains the lowest mean
SOH RMSE in the benchmark ($0.0292 \pm 0.0100$; the same backbone without the
physics loss scores $0.0377$) and the highest physical consistency of any
learned model ($0.884$); quantized to int8, an 11k-parameter deployment variant
measures 20.2~KB of flash, about 2\% of a Cortex-M4 budget. Negative results
(uncertainty weighting, cross-dataset extrapolation of the cycle coordinate,
the deployed-size accuracy trade-off) are reported with diagnosed mechanisms.

\section*{How the Phase-2 Feedback Has Been Addressed}

The examiner's feedback on the Phase-2 report and presentation is reproduced
below (verbatim, in italics), each item followed by the response and the
specific location in this thesis where it is integrated.

\subsection*{Report feedback}

\paragraph{R1.} \emph{``The figures must be improved.''}\\
\textbf{Response:} Every figure was regenerated from the reproducible pipeline
with publication settings: 600~dpi export, serif typography matching the thesis
body, a colorblind-safe palette with one fixed color per model across all
comparative figures, labeled axes and confidence intervals.
\textbf{Integration:} all figures in Chapters~\ref{ch:datasets_implementation}
and~\ref{ch:results}; new diagnostic figures include the differential-error
plot (Figure~\ref{fig:differential}), the $\alpha$-$\lambda$ prognosis
(Figure~\ref{fig:alpha_lambda}), and the footprint and accuracy-versus-memory
comparisons (Figures~\ref{fig:footprint} and~\ref{fig:frontier}).

\paragraph{R2.} \emph{``Some more explanation of ODE solving techniques, GRU and
motivation for other pre-processing methods should be provided.''}\\
\textbf{Response:} The ODE formulation now covers both laws, their integrated
closed forms, the sign convention of the rate constant, their role as soft
regularizers, and the precise definition of the derivative used in the physics
loss. The GRU section gives the full gate equations and the reasons for
choosing a GRU over an LSTM and a Transformer. Every preprocessing step
(Savitzky--Golay smoothing before ICA differentiation, per-cell imputation,
train-only MinMax scaling, windowing) is motivated where it is introduced.
\textbf{Integration:} Chapter~\ref{ch:problem}
(Sections~\ref{sec:physics_formulation}--\ref{sec:optimization_objective}),
Chapter~\ref{ch:methodology}
(Sections~\ref{sec:system_architecture}--\ref{sec:gru_component}),
Chapter~\ref{ch:datasets_implementation}
(Section~\ref{subsec:technical_breakdown}).

\paragraph{R3.} \emph{``Affiliation of mentor is missing in certificate
page.''}\\
\textbf{Response:} The certificate page has been corrected to state the
mentor's name together with his affiliation.
\textbf{Integration:} certificate page (front matter).

\subsection*{Presentation feedback}

\paragraph{P1.} \emph{``The presentation is fine. Mostly it was made by AI.''}\\
\textbf{Response:} The final thesis and its results are the candidate's own
work and are verifiable as such: every number, table and figure is generated by
a single reproducible pipeline from the raw datasets, interim placeholder
values have been replaced by measured results, and the corrections made during
the work --- including a data-leakage flaw in an input normalization that was
found and fixed during code review --- are documented inside the thesis rather
than hidden (Section~\ref{sec:caveats}). The manuscript has been revised
throughout for a uniform, formal academic style, and the candidate can
reproduce and explain the origin of every reported result.

\subsection*{Overall comments}

\paragraph{C1.} \emph{``The student should learn the fundamentals of the
problems and methodology.''}\\
\textbf{Response:} The final thesis demonstrates the fundamentals explicitly
where the Phase-2 draft was vague: the physics loss is defined precisely as the
partial derivative of the network output with respect to the cycle-coordinate
input and its role as a proxy for the trajectory derivative is stated
(Chapter~\ref{ch:problem}); the learned Verhulst parameters are reported with
their units and physical interpretation
(Section~\ref{sec:physics_interpretability}); the Physical Consistency Score is
bounded by the ground-truth reference so that it is read correctly as a
trend-plausibility measure (Section~\ref{sec:evaluation_strategy}); and the
statistical tests are performed at the correct unit of replication with their
limits stated (Section~\ref{sec:transformer_gap}).

\paragraph{C2.} \emph{``The student can completely avoid the ODE based approach
by integrating the logistic model and the double-exponential model.''}\\
\textbf{Response:} This suggestion was implemented in full as route~A: both
laws are used in their integrated closed forms, fit once on the training cells
and applied as a soft penalty, with no ODE solving. Rather than assuming one
mechanism is better, the thesis implements \emph{both} route~A and the
ODE-residual route~B and lets the data decide. The comparison is a central
result: the learnable route~B is better than or equal to the integrated fixed
prior on all four held-out cells and overall, so both the suggested approach
and the evidence for going beyond it are in the thesis.
\textbf{Integration:} Section~\ref{sec:two_routes}
(Chapter~\ref{ch:methodology}), Section~\ref{sec:central_finding} and
Table~\ref{tab:perfold} (Chapter~\ref{ch:results}).

\paragraph{C3.} \emph{``The student has to provide extensive baselines for
Physics informed machine learning methods for intended applications (extreme
hardware constraint).''}\\
\textbf{Response:} The benchmark was extended to ten models, all evaluated
under the same protocol: two non-neural floors (double-exponential curve fit,
SVR), LSTM, GRU, two Transformer scales, a continuous-time Neural-ODE as a
PIML-family comparator, a matched-capacity no-physics control, and the two
hybrid variants; the deployment candidates are additionally profiled for the
Cortex-M4 target. Broader PIML families are identified as future work.
\textbf{Integration:} Table~\ref{tab:benchmarks}
(Chapter~\ref{ch:datasets_implementation}), Tables~\ref{tab:benchmark}
and~\ref{tab:footprint} (Chapter~\ref{ch:results}), future work in
Chapter~\ref{ch:conclusion}.

\paragraph{C4.} \emph{``The student has to provide standard prognosis curves
for the example.''}\\
\textbf{Response:} Standard PHM prognostic evaluation was added: the
$\alpha$-$\lambda$ accuracy with the $\pm20\%$ cone and the Cumulative
Relative Accuracy, with RUL derived from the predicted SOH trajectory in the
standard way and post-EOL cycles masked so that only the informative pre-EOL
region is scored (the Prognostic Horizon degenerates under this masking and is
not used as a comparison axis).
\textbf{Integration:} Section~\ref{sec:phm}, Table~\ref{tab:phm} and
Figure~\ref{fig:alpha_lambda} (Chapter~\ref{ch:results}).

\paragraph{C5.} \emph{``Quantify the accuracy loss compared to the transformer
model; stating `85 to 90\%' is insufficient.''}\\
\textbf{Response:} The comparison is now fully quantified with a paired,
fold-level analysis against a high-capacity Transformer: the hybrid's mean SOH
RMSE is $21.0\%$ lower (though only 5 of 12 individual runs --- the mean gap
reflects the Transformer's run-to-run variance; not statistically significant
with four cells), and its physical consistency is higher on all 12 runs --- the
earlier qualitative claim is replaced by measured results with their
statistical limits stated.
\textbf{Integration:} Section~\ref{sec:transformer_gap} and
Table~\ref{tab:gap} (Chapter~\ref{ch:results}).

\paragraph{C6.} \emph{``Revise Figure 6.3 to clearly illustrate the deviation
in accuracy between models, potentially using a differential plot.''}\\
\textbf{Response:} A per-cycle differential-error plot was added exactly as
suggested: it shows the difference in absolute SOH error between each baseline
and the hybrid across the life of a held-out cell.
\textbf{Integration:} Figure~\ref{fig:differential}
(Chapter~\ref{ch:results}).

\paragraph{C7.} \emph{``Justify the aggregation of data points; articulate the
trade-offs and potential benefits beyond simply addressing time series data
requirements.''}\\
\textbf{Response:} The per-second to per-cycle aggregation is motivated in the
data chapter (a few hundred labeled cycles cannot support learning from raw
per-second streams; the ICA/DVA features retain the within-cycle information
that carries health signal), and the trade-off is quantified by a feature-set
ablation. The thesis states explicitly that this ablation is a feature-count
surrogate rather than a literal within-cycle representation study, which is
noted as future work.
\textbf{Integration:} Section~\ref{subsec:technical_breakdown}
(Chapter~\ref{ch:datasets_implementation}), Table~\ref{tab:aggregation}
(Chapter~\ref{ch:results}).

\paragraph{C8.} \emph{``Thoroughly investigate and present the rationale for
combining the [Verhulst] and double exponential ODEs, addressing concerns about
redundancy and potential for simpler solutions.''}\\
\textbf{Response:} The redundancy question was answered empirically with an
information-criterion analysis: the double-exponential attains a lower AIC than
the logistic on all five cells, so the two-timescale law is retained as the
primary prior and the logistic as an alternative configuration. The
full-protocol design ablation adds a complementary finding: used as soft
penalties inside the hybrid, the two priors perform near-identically --- the
decisive factor is not the prior's exact shape but whether its parameters are
learnable (route~B).
\textbf{Integration:} Section~\ref{sec:physics_interpretability},
Table~\ref{tab:prior_fits} and Figure~\ref{fig:redundancy}; ablation in
Table~\ref{tab:ablation} (Chapter~\ref{ch:results}).

\paragraph{C9.} \emph{``Demonstrate a deeper understanding of Bayesian
optimization and its applicability to the problem, particularly regarding the
weighting of data and physical loss.''}\\
\textbf{Response:} The thesis now separates the two distinct uses of
``Bayesian'' that were conflated earlier: Optuna's Tree-structured Parzen
Estimator for hyperparameter search, and the Bayesian (maximum-likelihood)
interpretation of homoscedastic uncertainty weighting for the data/physics loss
balance. The uncertainty weighting was then tested against a fixed weight on the full
benchmark protocol: it brings no reliable accuracy gain on either physics
route (on route~B its better mean is carried by run-to-run variance, 6 of 12
runs), with one consistent side effect --- a physical-consistency gain on
route~B. The comparison is reported in both directions and the fixed weight
is adopted for the headline models.
\textbf{Integration:} Sections~\ref{sec:uncertainty_weighting}--\ref{sec:two_routes}
(Chapter~\ref{ch:methodology}),
Table~\ref{tab:ablation} and Section~\ref{sec:ablations}
(Chapter~\ref{ch:results}).

\paragraph{C10.} \emph{``Explore and compare physics-informed machine learning
approaches as an alternative to the current hybrid model.''}\\
\textbf{Response:} Two PIML alternatives are compared: the continuous-time
Neural-ODE baseline, and, within the proposed framework, the fixed-prior
route~A against the learnable-residual route~B (two genuinely different
physics-informed mechanisms). The comparison of the two routes is the central
finding of the thesis; extending to further PIML families is listed as future
work.
\textbf{Integration:} Sections~\ref{sec:central_finding}
and~\ref{sec:main_benchmark} (Chapter~\ref{ch:results}),
Chapter~\ref{ch:conclusion}.

\paragraph{C11.} \emph{``Provide quantitative evidence of the memory footprint
reduction achieved with the hybrid model compared to transformers, specifically
in the context of target hardware (Cortex M4).''}\\
\textbf{Response:} The footprint is now measured, not estimated: the deployed
11k-parameter hybrid quantizes (dynamic int8) to 20.2~KB of flash with 1.25~KB
of activation RAM --- 2\% of the Cortex-M4 flash budget, thirty times less
flash than the high-capacity Transformer (611~KB) and $3.5\times$ less than the
edge-sized one (70~KB) --- with per-model MAC counts and an analytic latency
estimate labeled as such. A quantized ONNX export of the deployed model
accompanies the code.
\textbf{Integration:} Section~\ref{sec:resource_feasibility}
(Chapter~\ref{ch:datasets_implementation}), Section~\ref{sec:footprint} and
Table~\ref{tab:footprint} (Chapter~\ref{ch:results}).

\paragraph{C12.} \emph{``Revisit the definition of `ODE' and its application in
the context of the work, ensuring a clear understanding of its role and
benefits.''}\\
\textbf{Response:} Chapter~\ref{ch:problem} now defines both governing laws,
their integrated solutions, the meaning and sign of each parameter (including
why the fitted carrying capacity may exceed unity under rated-capacity
normalization and why the learned rate constant is negative), the soft-penalty
role of the ODE in the loss, and the exact derivative the implementation
computes. Chapter~\ref{ch:methodology} defines how the ODE enters each of the
two physics routes.
\textbf{Integration:} Chapter~\ref{ch:problem}
(Sections~\ref{sec:physics_formulation}--\ref{sec:optimization_objective}),
Section~\ref{sec:two_routes} (Chapter~\ref{ch:methodology}).

\section*{Summary}

All twelve overall comments and the three report-level comments have been
addressed in the final thesis; the presentation comment is answered by the
reproducibility of the work itself. Where a comment asked for a specific
artifact (differential plot, prognosis curves, Cortex-M4 quantification, AIC
redundancy analysis), that artifact now exists in Chapter~\ref{ch:results}.
Where a comment suggested a methodological change (integrated closed-form
physics instead of an ODE residual), the suggestion was implemented in full and
compared against the alternative, and the comparison itself became a
contribution of the thesis.
